\documentclass{article}
\usepackage{amsmath,amssymb}
\usepackage{xurl}
\usepackage[colorlinks=true,linkcolor=blue,citecolor=blue,urlcolor=blue]{hyperref}
% Shared commands for notes in docs/paper-gaps/.

\DeclareUnicodeCharacter{2080}{\ifmmode{}_{0}\else\textsubscript{0}\fi}
\DeclareUnicodeCharacter{2081}{\ifmmode{}_{1}\else\textsubscript{1}\fi}
\DeclareUnicodeCharacter{2082}{\ifmmode{}_{2}\else\textsubscript{2}\fi}
\DeclareUnicodeCharacter{2099}{\ifmmode{}_{n}\else\textsubscript{n}\fi}

\newcommand{\Lean}{\textsc{Lean}}
\ExplSyntaxOn
\NewDocumentCommand{\leanid}{m}{
  \ifmmode
    \textsf{\detokenize{#1}}
  \else
    \group_begin:
    \urlstyle{sf}
    \tl_set:Nn \l_tmpa_tl {#1}
    \tl_replace_all:Nnn \l_tmpa_tl {\_} {_}
    \exp_args:NV \path \l_tmpa_tl
    \group_end:
  \fi
}
\ExplSyntaxOff
\providecommand{\tr}{\operatorname{tr}}
\newcommand{\tnleanIssueBase}{https://github.com/LionSR/TNLean/issues/}
\newcommand{\tnleanPrBase}{https://github.com/LionSR/TNLean/pull/}
\newcommand{\tnleanDocBase}{https://sirui-lu.com/TNLean/docs/}
\newcommand{\ghissue}[1]{\href{\tnleanIssueBase#1}{GitHub issue \##1}}
\newcommand{\ghpr}[1]{\href{\tnleanPrBase#1}{PR \##1}}
\newcommand{\doclink}[1]{\href{\tnleanDocBase#1.html}{\path{#1}}}

% Dirac notation and the identity operator, as in blueprint/src/macros/common.tex.
% \providecommand keeps note-local definitions authoritative.
\usepackage{dsfont}
\providecommand{\ket}[1]{|#1\rangle}
\providecommand{\bra}[1]{\langle #1|}
\providecommand{\braket}[2]{\langle #1 | #2 \rangle}
\providecommand{\ketbra}[2]{|#1\rangle\!\langle #2|}
\providecommand{\Id}{\mathds{1}}
\providecommand{\one}{\mathds{1}}
\providecommand{\C}{\mathbb{C}}
\providecommand{\MN}[1]{M_{#1}(\C)}
\providecommand{\MD}{M_D(\C)}

% Verdict marker: \gapnote{<kind>}{<status>}. Kind is the policy
% classification (clarification, local-correction, scope-restriction,
% unfaithful, false-source, open-gap); status is open, wip (elimination
% underway), resolved, or historical. Machine-read by the site index;
% prints a verdict line.
\newcommand{\gapnote}[2]{\par\noindent\textbf{Verdict:} #1 (#2).\par}

\title{Rectangular and Vertex-Boundary Reduction to the Uniform Area Law}
\author{}
\date{October 17, 2026}
\begin{document}
\maketitle
\gapnote{scope-restriction}{resolved}

\section{The source statement}

Fix a local dimension $q\ge1$, an interaction norm bound $J>0$, and a spectral
gap $\Delta>0$. Corollary~1.2 of~\cite{OpenAI2026AreaLaw} states a uniform
entanglement-entropy bound for every subset of a finite rectangular lattice, in
both edge-boundary and vertex-boundary forms. Let
$V=\{0,\ldots,w-1\}\times\{0,\ldots,h-1\}$ be the rectangular grid of width $w$
and height $h$ with open nearest-neighbor edge set $E$. Consider a Hamiltonian
$H=\sum_{v\in V} h_v+\sum_{e\in E} h_e$, where each term $h_v$ or $h_e$ is
supported on its site $v$ or edge $e$ and has operator norm at most $J$. Only
Hermiticity of the sum $H=H^*$ is required, because replacing each term by
$(h_i+h_i^*)/2$ leaves $H$ unchanged without increasing the norm of any term.
Let $\Omega$ be a unit ground state of $H$ with ground energy $E_0$ satisfying
$H-E_0I\ge\Delta(I-|\Omega\rangle\langle\Omega|)$. For any region
$A\subseteq V$, let $S_\Omega(A)$ denote the von Neumann entropy of the reduced
state on $A$, let $\partial_V A$ be the set of edges in $E$ between $A$ and
$V\setminus A$, let $\partial_{\rm in}A$ be the set of vertices in $A$ adjacent
to $V\setminus A$, and let $Z_A$ be the set of vertices incident to
$\partial_V A$. Corollary~1.2 asserts that there is a constant
$C=C(q,1,J,\Delta)$, chosen before the side lengths $w,h$, the terms $h_v,h_e$,
the state $\Omega$, and the region $A$, such that
\begin{align}
 S_\Omega(A)&\le C|\partial_V A|,
 &S_\Omega(A)&\le4C|\partial_{\rm in}A|,
 &S_\Omega(A)&\le2C|Z_A|.
 \label{eq:rectangular_scope_bounds}
\end{align}
Empty boundaries are allowed, and for $q=1$ or $V=\emptyset$ every entropy
vanishes.\footnote{Source: \texttt{openai/math} revision
\texttt{adc7f1241b42e322a6451854ab7e4b4c146bf78a}, area-law
\texttt{00-introduction.tex}, lines 12--85, and
\texttt{01-preliminaries.tex}, lines 118--132. The total-Hermiticity convention
is also stated in the polynomial-PEPS manuscript, \texttt{00-introduction.tex},
lines 36--49.}

\section{The rectangular reduction}

Theorem~1.1 gives a uniform area law on arbitrary finite subsets of
$\mathbb Z^2$. The rectangular reduction applies this theorem to the grid $V$.
It embeds $V$ into $\mathbb Z^2$,
replaces each site and edge term by its Hermitian part on the corresponding
singleton or pair support in $\mathbb Z^2$, and sets the interaction to zero on
all other supports. Under the configuration isometry $U$ between the two
Hilbert spaces, the embedded Hamiltonian equals $UHU^*$. The proof establishes
that $U$ preserves the unit ground vector, the spectral gap $\Delta$, the
regional entropy $S_\Omega(A)$, and the cardinalities of $\partial_V A$,
$\partial_{\rm in}A$, and $Z_A$, and derives the vertex-boundary inequalities in
(\ref{eq:rectangular_scope_bounds}) from the degree bound of the rectangular
grid.\footnote{Declarations:
\leanid{TNLean.PEPS.AreaLaw.RectangularHamiltonian.toAreaLawHamiltonian},
\leanid{TNLean.PEPS.AreaLaw.RectangularHamiltonian.toAreaLawHamiltonian_operator},
\leanid{TNLean.PEPS.AreaLaw.RectangularHamiltonian.isGappedGroundState_toAreaLawHamiltonian},
\leanid{TNLean.PEPS.AreaLaw.regionalEntropy_rectangleAreaLawState},
\leanid{TNLean.PEPS.AreaLaw.card_edgeBoundary_rectangular},
\leanid{TNLean.PEPS.AreaLaw.card_innerBoundary_rectangular}, and
\leanid{TNLean.PEPS.AreaLaw.card_endpointBoundary_rectangular}.}

The uniform finite-domain area law is proved from the spectral-gap hypotheses.
Applying the rectangular reduction therefore proves Corollary~1.2 from the
lattice Hamiltonian hypotheses alone, with no area-law assumption.
The vertex-boundary inequalities also hold on arbitrary finite induced
domains.\footnote{Declarations:
\leanid{TNLean.PEPS.AreaLaw.uniformAreaLaw},
\leanid{TNLean.PEPS.AreaLaw.rectangularEntropyAreaLaw}, and
\leanid{TNLean.PEPS.AreaLaw.vertex_boundary_bounds}. The rectangular reduction
is \leanid{TNLean.PEPS.AreaLaw.UniformAreaLaw.rectangularEntropyAreaLaw}.}
Setting $w=h=L$ gives the square-grid area-law statement used by the
polynomial-PEPS construction, with its stated local dimension $q\ge2$ and no
further assumptions.\footnote{Declaration:
\leanid{TNLean.PEPS.AreaLaw.squareEntropyAreaLaw}.}

\section{Resolution of the restriction}

An earlier stage of the formalization took the uniform finite-domain area law
as an explicit assumption. The proof of Theorem~1.1 removes that assumption;
its geometric decomposition and collar amplification are established within
the argument. The resulting proofs of Corollary~1.2 and the square-grid area
law require no geometric or physical hypotheses beyond those in their source
statements. This note records the removal of an earlier restriction in the
formalization; it does not identify an error in the paper.

\bibliographystyle{alpha}
\bibliography{references}
\end{document}
